\section{Convergence Analysis}
To study convergence of a numerical solution we should define a measure for deviation between the analytical and numerical
approaches. In case of actual convergence we also need a factor that indicates speed of the process. Denote now
\begin{enumerate}
	\item \gls{u}: Exact solution of \cref{prb:Normalized_Problem}
	\item \gls{uNfem}: \Gls{FEMSol} of \cref{prb:Normalized_Problem} with $N$ elements. \\
				Note that $u_{N, \, fem}\left( x \right)$ might be
	\begin{itemize}
		\item \gls{uNlin}: Linear \gls{FEMSol} of \cref{prb:Normalized_Problem} with $N$ elements
		\item \gls{uNcub}: Cubic \gls{FEMSol} of \cref{prb:Normalized_Problem} with $N$ elements.
	\end{itemize}
\end{enumerate}
Define now following measures for solution deviations: \\
\begin{definition} \label{def:Squared_Error}
	In a Squared Error Measure we penalize the deviation between the numerical and analytical solutions by squaring its absolute
	value. Then we integrate this function over the whole region and take its square root, i.e.,
	\begin{equation}	\label{eqn:Squared_Error}
		\gls{ENsq} = \Bigg\{ \int_{x=0}^{1}\Big[ \gls{uEx} - \gls{uNfem} \Big] ^{2} \dd x \Bigg\}^{1/2}.
	\end{equation}
\end{definition}
%
\begin{definition}	\label{def:Max_Error}
	In a Max Error Measure we take absolute value of maximum deviation between the solutions, i.e.,
	\begin{align}	\label{eqn:Max_Error}
		\gls{ENmax} &= \max_{ 0 \leq x \leq 1 } \Big \vert  \gls{uEx} - \gls{uNfem} \Big \vert.
	\end{align}
\end{definition}
In this report we have chosen $E_{sq}$ for measuring the deviation. Having defined an error measure we move forward and define a
\gls{convFactor}.
%
\begin{definition} \label{def:Conv_Factor}
	Suppose that we have solved problem in \crefrange{eqn:FEM_Problem}{eqn:FEM_Right_Boundary} by applying \gls{FEM} with two different No.
	of elements, $N_{1}$ and $N_{2} = 2N_{1}$. Denote these solutions by $u_{1}$ and $u_{2}$ and let $E_{1}$
	and $E_{2}$ indicate the solution errors in the squared sense, see \cref{def:Squared_Error}. Define a \gls{convFactor} by
	\begin{equation} \label{eqn:Conv_Factor}
		\gls{fConv} = \frac{E_{1}}{E_{2}}.
	\end{equation}
\end{definition}
%
% Error estimates for the hyperbolic load
\begin{table}[tbh]
	\centering
	\rowcolors{1}{Gray}{LightCyan}
	\begin{tabular}{| c | c | c | c | c | c |} \hline
		% A constant
		\rowcolor{Cream} Load Function & $N$ & \gls{ENlin} & \gls{ENcub} & \gls{fConvLin} & \gls{fConvCub} \\ \hline
		$f_{n} \left( x \right) = 1$ & $4$ & \num{6.1 e-1} & \num{1.4 e-15} & - & -  \\ \hline
		$f_{n} \left( x \right) = 1$ & $8$	& \num{1.7 e-1} & \num{8.8 e-15} & $3.58$ & - \\ \hline
		$f_{n} \left( x \right) = 1$ & $16$ & \num{4.5 e-2} & \num{5.0 e-14} & $3.77$ & - \\ \hline
		% A linear function
		$f_{n} \left( x \right) = 1 + 2x$ & $4$	& \num{1.0 e0} & \num{5.6 e-4} & - & - \\ \hline
		$f_{n} \left( x \right) = 1 + 2x$ & $8$ & \num{2.8 e-1} & \num{3.4 e-5} & $3.61$&$16.3$ \\ \hline
		$f_{n} \left( x \right)=1 + 2x$ & $16$ & \num{7.3 e-2} & \num{2.1 e-6} & $3.79$ & $16.1$ \\ \hline
		% A quadratic function
		$f_{n} \left( x \right) = 1 + 2x - 3x^2$ & $4$ & \num{5.8 e-1} & \num{7.8 e-4} & - & - \\ \hline
		$f_{n} \left( x \right) = 1 + 2x - 3x^2$ & $8$ & \num{1.6 e-1} & \num{5.1 e-5} & $3.55$ & $15.3$ \\ \hline
		$f_{n} \left( x \right) = 1 + 2x - 3x^2$ & $16$	& \num{4.3 e-2} & \num{3.2 e-6} & $3.76$ & $15.8$ \\ \hline
	\end{tabular}
	\caption[Error Estimates for \glspl{FEMSol} of \cref{cas:q_1st_Degree_Denom}]
		{Error Estimates for \glspl{FEMSol} of \cref{cas:q_1st_Degree_Denom} with $e_1 = 10$ and $e_2 = 0.01$
		 and Boundary Parameters $\delta = 0$ and $P = 1$}
	\label{tbl:q_1st_Degree_Denom_Error}
\end{table}
%
% Error estimates for the classical problem with q(x) = e_1 / (x^2 + e_2)
\begin{table}
	\centering
	\rowcolors{1}{Gray}{LightCyan}
	\begin{tabular}{| c | c | c | c | c | c |} \hline
		\rowcolor{Cream} Load Function & $N$ & \gls{ENlin} & \gls{ENcub} & \gls{fConvLin} & \gls{fConvCub} \\ \hline
		% A constant
		$f_{n} \left( x \right) = 1$ & $4$ & \num{1.9 e-1} & \num{5.8 e-4} & - & - \\ \hline
		$f_{n} \left( x \right) = 1$ & $8$& \num{5.0 e-2} & \num{3.4 e-5} & $3.89$ & $17.0$	\\ \hline
		$f_{n} \left( x \right) = 1$ & $16$ & \num{1.3 e-2} & \num{2.1 e-6} & $3.94$ & $16.2$	\\ \hline
		% A linear function
		$f_{n} \left( x \right) = 1 + 2x$ & $4$	& \num{3.4 e-1} & \num{1.9 e-3} & - & - \\ \hline
		$f_{n} \left( x \right) = 1 + 2x$ & $8$	& \num{8.9 e-2} & \num{1.1 e-4} & $3.83$ & $16.7$ \\ \hline
		$f_{n} \left( x \right) = 1 + 2x$ & $16$ & \num{2.3 e-2} & \num{6.9 e-6} & $3.93$ & $16.2$ \\ \hline
		% A quadratic function
		$f_{n} \left( x \right) = 1 + 2x - 3x^2$ & $4$ & \num{1.7 e-1} & \num{8.6 e-4} & -& - \\ \hline
		$f_{n} \left( x \right) = 1 + 2x - 3x^2$ & $8$ & \num{4.2 e-2} & \num{6.2 e-5} & $3.94$ & $14.4$ \\ \hline
		$f_{n} \left( x \right) = 1 + 2x - 3x^2$ & $16$ & \num{1.1 e-2} & \num{4.0 e-6} & $3.95$ & $15.3$ \\ \hline
	\end{tabular}
	\caption[Error Estimates for \glspl{FEMSol} of \cref{cas:q_2nd_Degree_Denom}]
		{Error Estimates for \glspl{FEMSol} of \cref{cas:q_2nd_Degree_Denom} with $e_1 = 10$ and $e_2 = 0.01$ and Boundary Parameters $\delta = 0$ and $P = 1$}
	\label{tbl:q_2nd_Degree_Denom_Error}
\end{table}
%
% Error estimates for the exponential load
\begin{table}
	\centering
	\rowcolors{1}{Gray}{LightCyan}
	\begin{tabular}{| c | c | c | c | c | c |} \hline
		\rowcolor{Cream} Load Function & $N$ & \gls{ENlin} & \gls{ENcub} & \gls{fConvLin} & \gls{fConvCub} \\ \hline
		% A constant
		$f_{n} \left( x \right) = 1$ & $4$ & \num{1.3 e0}	& \num{4.2 e-3} & - & - \\ \hline
		$f_{n} \left( x \right) = 1$ & $8$	& \num{3.4 e-1} & \num{2.7 e-4} & $3.92$ & $15.6$ \\ \hline
		$f_{n} \left( x \right) = 1$ & $16$ & \num{8.5 e-2} & \num{1.7 e-5} & $3.98$ & $15.9$ \\ \hline
		% A linear function
		$f_{n} \left( x \right) = 1 + 2x$ & $4$	& \num{3.9 e0} & \num{1.9 e-2} & - & - \\ \hline
		$f_{n} \left( x \right) = 1 + 2x$ & $8$	& \num{1.0 e0} & \num{1.2 e-3} & $3.87$ & $15.5$ \\ \hline
		$f_{n} \left( x \right) = 1 + 2x$ & $16$ & \num{2.5 e-1} & \num{7.8 e-5} & $3.97$ & $15.9$ \\ \hline
		% A quadratic function
		$f_{n} \left( x \right) = 1+2x-3x^2$ & $4$ & \num{1.1 e0} & \num{1.7 e-2} & - & - \\ \hline
		$f_{n} \left( x \right) = 1+2x-3x^2$ & $8$ & \num{2.9 e-1} & \num{1.1 e-3} & $3.87$ & $15.0$ \\ \hline
		$f_{n} \left( x \right) = 1+2x-3x^2$ & $16$	& \num{7.4 e-2} & \num{7.0 e-5} & $3.96$ & $15.7$ \\ \hline

	\end{tabular}
	\caption[Error Estimates for \glspl{FEMSol} of \cref{cas:q_Exponential}]
		{Error Estimates for \glspl{FEMSol} of \cref{cas:q_Exponential} with $\alpha = 1$ and Boundary Parameters $\delta = 0$ and $P = 0.01$}
	\label{tbl:q_Exponential_Error}
\end{table}
Then the solution is convergent if $f_{con} > 1$ for all choices of $N_{1}$. Furthermore the larger $f_{con}$, the faster convergence. A value of $f_{c} = 4$, for example, indicates that doubling No. of elements results in a
reduction of the error by a factor $4$. Speed of convergence in this case is said to be of $\mathcal{O}\left( N^2\right) $. \\

Convergence speed of \glspl{FEMSol} in \cref{cas:q_1st_Degree_Denom,cas:q_2nd_Degree_Denom,cas:q_Exponential} is shown in \crefrange{tbl:q_1st_Degree_Denom_Error}{tbl:q_Exponential_Error}. A closer look at \cref{fig:q_1_Load_0} and \cref{tbl:q_1st_Degree_Denom_Error,tbl:q_2nd_Degree_Denom_Error} reveals the following facts:
\begin{enumerate}
	\item For a given \gls{N} it seems that both linear and cubic \glspl{FEMSol} stay on the same side of the \gls{analSol}, i.e., they do not oscillate around the \gls{analSol}.

	\item \Gls{linearFEMSol} seems to underestimate the displacement solution. \Gls{cubicFEMSol} on the other hand, changes between being larger and smaller than the \gls{analSol}.

	\item \Glsentryname{order0Cont} in \gls{linearFEMSol} is clearly seen when No. of elements is small, look at the \glspl{nodePoint} in upper left part of \cref{fig:q_1_Load_0}. Cubic \glspl{FEMSol} on the other hand,
	has continuous first derivative, see lower parts of \cref{fig:q_1_Load_0}.

	\item Obviously we can not satisfy the \gls{naturalCondition} with the \gls{linearFEMSol}, because it involves value of first derivative. In the cubic case, on the other hand, the solution has the exact derivative on the \gls{rightBoundary}.

	\item The linear and cubic \glspl{FEMSol} seem to converge as $\mathcal{O} \left( h^2 \right)$ and $\mathcal{O} \left( h^4 \right)$. This is in accordance with the results in \cref{cha:Error_Estimates}, see \cref{eqn:FEM_Error_To_h_And_Exact_Sol}.

	\item \Cref{fig:q_1_Load_0,tbl:q_1st_Degree_Denom_Error} suggest a cubic FEM-error comparable to the numerical round off error. This should not come as a surprise because the problem's \gls{analSol} is also a third degree polynomial, see \cref{eqn:q_1st_Degree_Denom_Solution}.

	\item Convergence speed decreases as by growing \gls{N}. This is of course due the round off error. This fact is more clear for cubic \gls{FEMSol} because of its faster speed of convergence.
\end{enumerate}	\text{}
%
Finally, the code in \cref{lst:Estimate_Error} estimates the squared error of \glspl{FEMSol} compared to the \gls{analSol} in each case.
